% ======================================================================
% GENERATED by tools/build-language-preambles.py from
%   grundstrukturen-ai-transcription-spring-2025\grundstrukturen.tex
% Do not edit by hand: edit the book preamble and rerun the script.
% ======================================================================

% =====================================================================
%  SHORT-NAME ENVIRONMENT VOCABULARY
%  Converted in place from the prompt repo's tools/make-preamble.py.
%
%    spoken-clean         -> speech        math-stroke          -> content
%    student-interaction  -> student       explanation-of-steps -> steps
%    didactic-insight     -> insight       meta-note            -> note
%    nice-box             -> overview      color-box            -> highlight
%    lecture-break        -> pause
%    short-proof          -> proof         proof                -> proofWrapper
%    ai-note              -> aiNote
%    ai-*-invisible-content -> hidden
% =====================================================================

\documentclass[11pt,a4paper]{report}

% ==========================================
% NEW VIBRANT SHADOW PREAMBLE
% ==========================================

% pdfLaTeX (grundstrukturen.pdf) sets the text in 8-bit T1 fonts. LuaLaTeX (grundstrukturen-v1.pdf)
% needs Unicode fonts: it hands every character to the font as it is, and in T1 slot 223 (the
% code of ß) holds "SS", while characters such as • and — do not exist there at all.
\usepackage{iftex}
\ifPDFTeX
  \usepackage[utf8]{inputenc}
  \usepackage[T1]{fontenc}
\fi

% --- FONT CONFIGURATION ---
\ifPDFTeX
  \usepackage{lmodern}
  \usepackage{anyfontsize}
  \usepackage{mathpazo}
  \renewcommand{\sfdefault}{lmss}
\else
  \usepackage{anyfontsize}
  \usepackage{mathpazo}                 % math as under pdfLaTeX
  \usepackage[no-math]{fontspec}
  \setmainfont{TeX Gyre Pagella}        % Palatino, like mathpazo's text font
  \setsansfont{Latin Modern Sans}
  \setmonofont{Latin Modern Mono}
\fi

% Schweizer Rechtschreibung wie im Produktions-Repo: ß wird beim Setzen zu ss (heisst, gross),
% Quelltext und Label-Schlüssel behalten ß. nswissgerman statt swissgerman, damit nach der neuen
% Rechtschreibung getrennt wird (Kis-te, nicht Ki-ste).
\usepackage[nswissgerman]{babel}
\usepackage{newunicodechar}
\newunicodechar{ß}{ss}
\usepackage{csquotes}

% --- OPTIMIERTE SILBENTRENNUNG (HYPHENATION TUNING) ---
\hyphenpenalty=100
\exhyphenpenalty=100
\doublehyphendemerits=10000
\finalhyphendemerits=5000
\tolerance=1500
\hbadness=1500
\emergencystretch=3em
\hfuzz=1pt

\usepackage{amsmath, amssymb, amsthm, mathtools, environ, etoolbox} 
\usepackage{booktabs}
\usepackage{array}
\newcolumntype{L}[1]{>{\raggedright\arraybackslash}p{#1}}
\usepackage{listings}
\usepackage{capt-of}
\usepackage{fontawesome5}
\usepackage[dvipsnames, x11names]{xcolor}
\usepackage{siunitx}
\usepackage[version=4]{mhchem}
\usepackage{bm}
\usepackage{mathrsfs}
\usepackage{cancel}
\usepackage{braket}
\usepackage{derivative}
\usepackage{xfrac, faktor}            % \sfrac{a}{b}, \faktor{R}{K}: slanted fractions and quotients
\usepackage{dsfont}                   % \mathds{1}, the indicator function or identity
\usepackage[normalem]{ulem}           % \sout{...}, a word struck out on the board (normalem keeps \emph)
\usepackage{polynom}                  % \polylongdiv, polynomial long division
\usepackage{tabularx, multirow, diagbox, colortbl} % tables: \multirow, \diagbox corner cell (Cayley tables), \cellcolor
\usepackage{nicematrix}               % matrices with dotted lines, augmented-matrix bars, row-operation labels
\usepackage[only, llbracket, rrbracket, mapsfrom, longmapsfrom, lightning]{stmaryrd} % only these, the rest would override amssymb; \lightning is the board's contradiction mark
\usepackage{algorithm, algpseudocode}

% Setzt global "Woche" anstelle von "Kapitel" (angepasst an nswissgerman)
\addto\captionsnswissgerman{\renewcommand{\chaptername}{Woche}}
\addto\captionsnswissgerman{\renewcommand{\contentsname}{Inhaltsverzeichnis}}

% ==========================================
% GLOBAL COLOR PALETTE & VARIABLES
% ==========================================
% 0. CUSTOM COLORS
\definecolor{DeepSeaGreen}{HTML}{4A6F6E}
\definecolor{DeepSeaGreenBright1}{HTML}{335655}
\definecolor{DeepSeaGreenBright2}{HTML}{4A6F6E}
\definecolor{DarkCaramel}{HTML}{684A31}

% 1. CORE THEME COLORS
\colorlet{ThemePrimary}{RoyalPurple}
\colorlet{ThemeSecondary}{MidnightBlue}
\colorlet{ThemeWeekNumber}{OliveGreen}
\colorlet{ThemeAccent1}{BurntOrange}
\colorlet{ThemeAccent2}{TealBlue}
\colorlet{ThemeAccent3}{ForestGreen}
\colorlet{ThemeAccent4}{BrickRed}

% 2. TEXT & LINKS
\definecolor{TextEmphColor}{HTML}{BD6400}
\definecolor{TextBoldColor}{HTML}{0E6C95}
\colorlet{TextMetaNote}{Gray}
\colorlet{LinkColor}{MidnightBlue}
\colorlet{FileColor}{Magenta}
\definecolor{ListAmber}{HTML}{BD2400}
\colorlet{ItemizeBullet}{ListAmber}
\colorlet{EnumerateBullet}{ListAmber}
\definecolor{QuoteRed}{HTML}{BD2400}
\colorlet{HeadingSlate}{SlateBlue4}
\colorlet{HeadingBrown}{Brown!80!black}

% --- TOC ENTRY COLOURS ---
\definecolor{TocEntryBlue}{HTML}{008CC8}
\definecolor{TocEntryBrown}{HTML}{A8651E}
\colorlet{TocPartColor}{ThemePrimary!90!black}
\colorlet{TocChapColor}{ThemeSecondary!90!black}
\colorlet{TocSecColor}{TocEntryBlue}
\colorlet{TocSubsecColor}{TocEntryBrown}
\colorlet{TocSubsubColor}{gray}
\colorlet{TocLabelColor}{ThemePrimary!85!black}
\colorlet{TocProofColor}{OliveGreen}

% 3. SHADOWS & PROOF GLOBALS
\colorlet{ShadowNormal}{black!50}
\colorlet{ShadowProof}{OliveGreen!40}
\colorlet{ProofTint}{OliveGreen}

% 4. ENVIRONMENTS BASE COLORS
\colorlet{EnvStudent}{Violet}
\colorlet{EnvSpoken}{MidnightBlue}
\colorlet{EnvMath}{MidnightBlue}
\colorlet{EnvTeacher}{ForestGreen}
\colorlet{EnvMeta}{TealBlue}
\colorlet{EnvAI}{BrickRed}
\colorlet{EnvNote}{BurntOrange}
\colorlet{EnvBreak}{BrickRed}

% --- THEOREM ENVIRONMENT COLORS ---
\definecolor{OriginalTheoremHex}{HTML}{6A3199}
\colorlet{EnvTheorem}{OriginalTheoremHex}
\colorlet{EnvTheoremDark}{OriginalTheoremHex}
\definecolor{EnvRemark}{HTML}{008837}           % Bemerkung, Notation, Zusammenfassung, Ziele
\colorlet{EnvQuestion}{MidnightBlue}            % Frage, Antwort
\colorlet{EnvQuestion}{EnvRemark}            % pipeline: questions/answers green like \recap
\definecolor{EnvImportantRemark}{HTML}{D7263D}  % Wichtige Bemerkung

% --- EMPHASIS STYLE OVERRIDE ---
% \emph and \textbf in sans-serif bold, each in its color, as in the other books.
\DeclareTextFontCommand{\emph}{\sffamily\bfseries\color{TextEmphColor}}
\DeclareTextFontCommand{\textbf}{\sffamily\bfseries\color{TextBoldColor}}

% 5. ENVIRONMENT MIXES & PROOF STATES
\colorlet{EnvSpokenBg}{EnvSpoken!5}
\colorlet{EnvSpokenProof}{EnvSpoken!50!ProofTint}
\colorlet{EnvNoteBg}{EnvNote!20}
\colorlet{EnvNoteProof}{EnvNote!50!ProofTint}
\colorlet{EnvTeacherBg}{EnvTeacher!5}
\colorlet{EnvTeacherBorder}{EnvTeacher!80!black}
\colorlet{EnvTeacherProof}{EnvTeacherBorder!50!ProofTint}
\colorlet{EnvMathProof}{EnvMath!50!ProofTint}
\colorlet{ProofSubBg}{EnvMath!30}
\colorlet{ProofSubText}{EnvMath!80!black}
\colorlet{EnvMetaProof}{EnvMeta!50!ProofTint}
\colorlet{EnvAIBg}{EnvAI!5}
\colorlet{EnvAIBorder}{EnvAI!85!black}
\colorlet{EnvAIProof}{EnvAIBorder!50!ProofTint}
\colorlet{EnvAIFallbackProof}{EnvAI!50!ProofTint}
\colorlet{EnvStudentBg}{EnvStudent!5}
\colorlet{EnvStudentProof}{EnvStudent!50!ProofTint}
\colorlet{EnvBreakProof}{EnvBreak!50!ProofTint}
\colorlet{ExplanationIconColor}{ProofTint}

% 6. PROOF SPECIFIC ELEMENTS
\colorlet{ProofQEDText}{ProofTint}
\colorlet{ProofQEDBg}{white}

% 7. HEADER, FOOTER & CHAPTER
\colorlet{HeaderFooterFaint}{ThemeWeekNumber!45!Gray}
\colorlet{HeaderFooterLine}{HeaderFooterFaint!50}
\colorlet{ChapTitleText}{ThemeSecondary!90!black}
\colorlet{ChapNumberText}{ThemeWeekNumber}
\colorlet{ChapSubtitleText}{HeaderFooterFaint}
\colorlet{ChapLine}{ThemeSecondary!80!black}

% 8. SECTION BARS
\colorlet{SecBar1}{ThemePrimary}
\colorlet{SecBar2}{HeadingBrown}
\colorlet{SecBar3}{ThemeAccent3!80!black}
\colorlet{SecBar4}{ThemeAccent4!90!black}
\colorlet{SecBar1Proof}{SecBar1!50!ProofTint}
\colorlet{SecBar2Proof}{SecBar2!50!ProofTint}
\colorlet{SecBar3Proof}{SecBar3!50!ProofTint}
\colorlet{SecBar4Proof}{SecBar4!50!ProofTint}
\colorlet{SecTitleColor}{ThemePrimary}
\colorlet{SecTitleColorProof}{SecTitleColor!50!ProofTint}

\colorlet{SubSecBar1}{HeadingBrown}
\colorlet{SubSecBar2}{HeadingBrown!70!white}
\colorlet{SubSecBar1Proof}{SubSecBar1!50!ProofTint}
\colorlet{SubSecBar2Proof}{SubSecBar2!50!ProofTint}
\colorlet{SubSecTitleColor}{HeadingBrown}
\colorlet{SubSecTitleColorProof}{SubSecTitleColor!50!ProofTint}

\colorlet{SubSubSecBar1}{HeadingBrown}
\colorlet{SubSubSecBar2}{gray!60!white}
\colorlet{SubSubSecBar1Proof}{SubSubSecBar1!50!ProofTint}
\colorlet{SubSubSecBar2Proof}{SubSubSecBar2!50!ProofTint}
\colorlet{SubSubSecTitleColor}{HeadingBrown}
\colorlet{SubSubSecTitleColorProof}{SubSubSecTitleColor!50!ProofTint}
\colorlet{SubSubSecNumberColor}{gray}
\colorlet{SubSubSecNumberColorProof}{SubSubSecNumberColor!50!ProofTint}

% ==========================================
% PAGE GEOMETRY & HYPERREF
% ==========================================
\usepackage[
  a4paper,
  top=24mm,
  bottom=26mm,
  inner=26mm,
  outer=26mm,
  headheight=14pt,  % one-line running heads (see HEADER, FOOTER below)
  headsep=7mm,
  footskip=11mm
]{geometry}

\usepackage{bookmark}
\usepackage{hyperref}
\hypersetup{
  colorlinks=true,
  linkcolor=LinkColor,
  filecolor=FileColor,
  urlcolor=LinkColor,
  citecolor=LinkColor,
  pdftitle={Grundstrukturen -- Vorlesungstranskript},
  pdfauthor={Christian Urech / ETH Z{\"u}rich},
  pdfsubject={Grundstrukturen},
  bookmarksnumbered=true,
  bookmarksopen=true
}
\usepackage[nameinlink,capitalize,noabbrev]{cleveref}

% ==========================================
% BOXES, SHADOWS & ENGINE SETUP
% ==========================================
\usepackage{varwidth}
\usepackage[most]{tcolorbox}
\tcbuselibrary{skins, breakable, theorems}

\newif\ifinsideproof
\insideprooffalse

\tcbset{
  dynamic shadow/.code={
    \ifinsideproof
      \tcbset{drop fuzzy shadow=ShadowProof} 
    \else
      \tcbset{drop fuzzy shadow=ShadowNormal}
    \fi
  },
  dynamic color shadow/.code={
    \ifinsideproof
      \tcbset{drop fuzzy shadow=ShadowProof} 
    \else
      \tcbset{drop fuzzy shadow=#1}
    \fi
  }
}

\tcbsetforeverylayer{
  before skip=14pt plus 3pt minus 2pt, 
  after skip=14pt plus 3pt minus 2pt,
  pad at break=2.5mm,
  lines before break=2,
  shrink break goal=2mm,
  ignore nobreak=false,
  % \textbf in a box title takes the title's color: blue bold on a red title bar is unreadable
  before title app={\colorlet{TextBoldColor}{.}},
  % No shadow at a page break: the parts of a broken box get their shadow along the cut edges
  % too, so white patches cover it there (underlays are drawn after the shadow).
  underlay first={%
    \fill[white] ([xshift=-15pt]frame.south west) rectangle ([xshift=15pt,yshift=-15pt]frame.south east);},
  underlay middle={%
    \fill[white] ([xshift=-15pt]frame.south west) rectangle ([xshift=15pt,yshift=-15pt]frame.south east);
    \fill[white] ([xshift=-15pt]frame.north west) rectangle ([xshift=15pt,yshift=15pt]frame.north east);},
  underlay last={%
    \fill[white] ([xshift=-15pt]frame.north west) rectangle ([xshift=15pt,yshift=15pt]frame.north east);}
}

% ==========================================
% TYPOGRAPHY & SPACING
% ==========================================
\linespread{1.05}

\newlength{\standardparindent}
\setlength{\standardparindent}{17pt} 

\newlength{\customparskip}
\setlength{\customparskip}{3pt plus 1.5pt minus 0.5pt}

\setlength{\parindent}{0pt}
\setlength{\parskip}{\customparskip}

\newlength{\customenvspace}
\setlength{\customenvspace}{0.9ex plus 0.2ex minus 0.1ex}

\setlength{\mathsurround}{1pt}
\renewcommand{\arraystretch}{1.15}
\setlength{\intextsep}{1.5\customparskip plus 1ex minus 1ex}
\setlength{\textfloatsep}{2\customparskip plus 1ex minus 1ex}

\AtBeginDocument{
  \setlength{\abovedisplayskip}{1.2\customparskip plus 0.5ex minus 0.2ex}
  \setlength{\belowdisplayskip}{1.2\customparskip plus 0.5ex minus 0.2ex}
  \setlength{\abovedisplayshortskip}{0.5\customparskip plus 0.2ex}
  \setlength{\belowdisplayshortskip}{0.8\customparskip plus 0.2ex minus 0.1ex}
}

\usepackage{enumitem}
\setlist{
  topsep=0pt, 
  partopsep=0pt, 
  parsep=0.5\customparskip, 
  itemsep=0pt, 
}
\setlist[itemize]{label=\textcolor{ItemizeBullet}{\textbullet}}
\setlist[enumerate]{
  % Every label in the enumerate color and sans-serif bold, however it is written: the default
  % (a), label=(1), or label=\textbf{(i)} (TextBoldColor is redirected for \textbf).
  font=\sffamily\bfseries\color{EnumerateBullet}\colorlet{TextBoldColor}{EnumerateBullet},
  label=(\alph*),
  ref=(\alph*)
}
\setlist[description]{
  style=unboxed,
  labelsep=0.333em plus 0.167em minus 0.111em
}
\setlist[itemize,2]{topsep=0pt, itemsep=0pt, parsep=0pt}
\setlist[enumerate,2]{topsep=0pt, itemsep=0pt, parsep=0pt}

\makeatletter
\newcommand{\fixlistaftertheorem}{%
  \if@inlabel 
    \leavevmode \par 
  \fi
}
\BeforeBeginEnvironment{itemize}{\fixlistaftertheorem}
\BeforeBeginEnvironment{enumerate}{\fixlistaftertheorem}
\BeforeBeginEnvironment{description}{\fixlistaftertheorem}
\makeatother

% ==========================================
% HEADER, FOOTER & COMMANDS
% ==========================================
\usepackage{fancyhdr}
\usepackage{lastpage}
\usepackage[fit]{truncate}
\pagestyle{fancy}
\fancyhf{}

% Running heads, one line each: the header carries the section (left) and week, date and page
% (right); the footer the subsection (left) and subsubsection (right). Each heading level has its
% own mark class, and \LastMark gives the one in force at the bottom of the page. A title too long
% for its slot is cut at a word and ends in "...", so a head never wraps to a second line. Without
% a subsubsection the subsection gets the whole footer.
\NewMarkClass{gs-sec}
\NewMarkClass{gs-subsec}
\NewMarkClass{gs-subsubsec}
\newcommand{\headfont}{\footnotesize\sffamily\color{HeaderFooterFaint}}
\newcommand{\pagecounter}{\textbf{\thepage} / \pageref*{LastPage}}

\fancyhead[L]{\headfont\truncate{0.6\headwidth}{\LastMark{gs-sec}}}
\fancyhead[R]{\headfont\leftmark\quad\pagecounter}

\fancyfoot[L]{\headfont\sbox0{\LastMark{gs-subsubsec}}%
  \truncate{\ifdim\wd0=0pt \headwidth\else 0.49\headwidth\fi}{\LastMark{gs-subsec}}}
\fancyfoot[R]{\headfont\truncate{0.49\headwidth}{\LastMark{gs-subsubsec}}}

% A plain \hrule, not an \hbox: a box in the head's vbox gets a full \baselineskip above it,
% which pushed the line out of the one-line head (it vanished, and the head became 21pt tall).
\renewcommand{\headrule}{\vskip3pt{\color{HeaderFooterLine}\hrule height 0.5pt width \headwidth}\vskip-3.5pt}
\renewcommand{\footrule}{\hbox to\headwidth{\color{HeaderFooterLine}\leaders\hrule height 0.5pt\hfill}}

\fancypagestyle{plain}{%
  \fancyhf{}%
  \fancyfoot[C]{\headfont\pagecounter}%
  \renewcommand{\headrule}{}%
  \renewcommand{\footrule}{}%
}

% A new heading clears the levels below it, so a page never shows a subsection of the last section.
\renewcommand{\chaptermark}[1]{}
\renewcommand{\sectionmark}[1]{%
  \InsertMark{gs-sec}{\thesection.\ #1}%
  \InsertMark{gs-subsec}{}\InsertMark{gs-subsubsec}{}%
}
\renewcommand{\subsectionmark}[1]{%
  \InsertMark{gs-subsec}{\thesubsection.\ #1}\InsertMark{gs-subsubsec}{}%
}
\renewcommand{\subsubsectionmark}[1]{%
  \InsertMark{gs-subsubsec}{\thesubsubsection.\ #1}%
}

% TikZ Configuration
\usepackage{tikz}
\usetikzlibrary{calc, arrows.meta, positioning, 3d, shapes.geometric, decorations.pathreplacing, patterns, shadows}
\usetikzlibrary{fit}
\usetikzlibrary{arrows}
\usetikzlibrary{babel}
\usetikzlibrary{angles, quotes, intersections}
\usetikzlibrary{automata, matrix, backgrounds, shapes.misc, shapes.arrows, decorations.markings, decorations.pathmorphing}
\usepackage{tikz-cd}
\usepackage{pgfplots}
\pgfplotsset{compat=1.18}
\usepgfplotslibrary{fillbetween}
\usepackage{circuitikz}
\usepackage[nobottomtitles*]{titlesec}
\renewcommand{\bottomtitlespace}{3.5\baselineskip}
\usepackage{titletoc}
\usepackage{needspace}  % keeps a "Woche N" TOC heading with its first entries

\setcounter{secnumdepth}{3}
\setcounter{tocdepth}{3}

% ==========================================
% TABLE OF CONTENTS TYPOGRAPHY & COLORS
% ==========================================
\contentsmargin{4.7ex}

\newcommand{\tocindent}[1]{\ifcase\numexpr#1\relax 0pt\or 4.2ex\or 8.5ex\or 13.5ex\or 18.5ex\or 23.5ex\else 28.5ex\fi}

\makeatletter
\def\proofwrapper@tocparent{-1}
\def\proofwrapper@tocshift{-1}
\newcommand{\proofwrappertocbegin}[1]{\gdef\proofwrapper@tocparent{#1}\gdef\proofwrapper@tocshift{-1}}
\newcommand{\proofwrappertocend}{\gdef\proofwrapper@tocparent{-1}}
\newcommand{\proofwrapper@tocinner}[1]{%
  \ifnum\proofwrapper@tocparent<\z@\else
    \ifnum\proofwrapper@tocshift<\z@
      \count@=\numexpr\proofwrapper@tocparent+2-#1\relax
      \ifnum\count@<\z@ \count@=\z@ \fi
      \xdef\proofwrapper@tocshift{\the\count@}%
    \fi
    \setlength\leftskip{\tocindent{#1+\proofwrapper@tocshift}}%
  \fi}
% titletoc needs a level for the entry type: shown up to tocdepth 3
\def\ttll@proofwrapper{3}

% Level 0 (Part): Weekly Blocks
\titlecontents{part}[0ex]
  {\needspace{6\baselineskip}\addvspace{3.5ex}\hypersetup{linkcolor=TocPartColor}\sffamily\bfseries\large\color{TocPartColor}}
  {}
  {}
  {}
  [\vspace{0.8ex}{\color{HeaderFooterLine}\titlerule[0.8pt]}\vspace{1.2ex}]

% Level 1 (Chapter): Individual Lecture
\newcommand{\tocchaptericon}{\faCalendarDay}
\titlecontents{chapter}[0ex]
  {\addvspace{2.2ex}\hypersetup{linkcolor=TocChapColor}\sffamily\bfseries\normalsize\color{TocChapColor}}
  {\color{ThemePrimary!85!black}\tocchaptericon\hspace{0.6em}}
  {}
  {\color{HeaderFooterLine}\titlerule*[0.6pc]{.}\color{ThemeWeekNumber}\contentspage}

% Level 2 (Section): Lecture Sections
\titlecontents{section}[\tocindent{1}]
  {\addvspace{0.8ex}\proofwrapper@tocinner{1}\hypersetup{linkcolor=TocSecColor}\sffamily\small\color{TocSecColor}}
  {\color{TocLabelColor}\thecontentslabel\hspace{1.2ex}\color{TocSecColor}}
  {}
  {\color{HeaderFooterLine}\titlerule*[0.6pc]{.}\color{TocSecColor}\contentspage}

% Level 3 (Subsection): Subsections
\titlecontents{subsection}[\tocindent{2}]
  {\addvspace{0.4ex}\proofwrapper@tocinner{2}\hypersetup{linkcolor=TocSubsecColor}\sffamily\footnotesize\color{TocSubsecColor}}
  {\color{TocSubsecColor}\thecontentslabel\hspace{1.2ex}\color{TocSubsecColor}}
  {}
  {\color{HeaderFooterLine}\titlerule*[0.6pc]{.}\color{TocSubsecColor}\contentspage}

% Level 4 (Subsubsection): Subsubsections
\titlecontents{subsubsection}[\tocindent{3}]
  {\addvspace{0.2ex}\proofwrapper@tocinner{3}\hypersetup{linkcolor=TocSubsubColor}\sffamily\footnotesize\color{TocSubsubColor}}
  {\color{TocSubsubColor}\thecontentslabel\hspace{1.2ex}\color{TocSubsubColor}}
  {}
  {\normalfont\color{HeaderFooterLine}\titlerule*[0.6pc]{.}\color{TocSubsubColor}\contentspage}

% Titled proofWrapper
\titlecontents{proofwrapper}[0pt]
  {\addvspace{0.2ex}\setlength\leftskip{\tocindent{\proofwrapper@tocparent+1}}\hypersetup{linkcolor=TocProofColor}\sffamily\footnotesize\color{TocProofColor}}
  {}
  {\faPenNib\hspace{0.6em}\itshape}
  {\normalfont\sffamily\color{HeaderFooterLine}\titlerule*[0.6pc]{.}\color{TocProofColor}\contentspage}
\makeatother

% Math operators
\DeclareMathOperator{\id}{id}
\DeclareMathOperator{\Vol}{Vol}
\providecommand{\bigtimes}{\mathop{\mathchoice
  {\vcenter{\hbox{\scalebox{1.8}{$\displaystyle\times$}}}}
  {\vcenter{\hbox{\scalebox{1.35}{$\textstyle\times$}}}}
  {\vcenter{\hbox{\scalebox{1.2}{$\scriptstyle\times$}}}}
  {\vcenter{\hbox{\scalebox{1.2}{$\scriptscriptstyle\times$}}}}}\displaylimits}
\providecommand{\dv}{\odv}
% Operators and delimiters Gemini commonly writes; transcription.md lists them as preloaded.
\providecommand{\rank}{\operatorname{rank}}
\providecommand{\Hom}{\operatorname{Hom}}
\providecommand{\End}{\operatorname{End}}
\providecommand{\Aut}{\operatorname{Aut}}
\providecommand{\sgn}{\operatorname{sgn}}
\providecommand{\diag}{\operatorname{diag}}
\providecommand{\tr}{\operatorname{tr}}
\providecommand{\Span}{\operatorname{span}}
\providecommand{\im}{\operatorname{im}}
\providecommand{\lcm}{\operatorname{lcm}}
\ifdefined\abs\else\DeclarePairedDelimiter\abs{\lvert}{\rvert}\fi   % \abs{x}, \abs*{x} grows
\ifdefined\norm\else\DeclarePairedDelimiter\norm{\lVert}{\rVert}\fi % \norm{v}

% Custom text commands (Deutsch)
\newcommand{\qt}[1]{\textit{\textcolor{QuoteRed}{\glqq}#1\textcolor{QuoteRed}{\grqq}}}
\newcommand{\newterm}[1]{\textcolor{QuoteRed}{\glqq}\textcolor{TextBoldColor}{\textit{#1}}\textcolor{QuoteRed}{\grqq}}
\newcommand{\inlinemetanote}[1]{{\colorlet{QuoteRed}{TextMetaNote}\textit{\textcolor{TextMetaNote}{[#1]}}}} % quote marks and brackets grey too
\newcommand{\recap}[1]{\par\smallskip\noindent{\sffamily\bfseries\color{EnvRemark}\faHistory\ R{\"u}ckblick:}\ #1\par\smallskip}
\NewDocumentCommand{\checkline}{ O{Probe} m }{\par\smallskip\noindent{\sffamily\bfseries\color{EnvRemark}\faCheckCircle\ #1:}\ #2\par\smallskip}

% Custom Theorem Styles (Deutsch)
\newcommand{\NewIconThmStyle}[3]{%
\newtheoremstyle{#1}%
{\customenvspace}{\customenvspace}{\normalfont}{}{\color{#2}\sffamily\bfseries}{:}{0.5em}%
{#3\hspace{0.4em}\thmname{##1}\thmnumber{ ##2}\thmnote{ \textcolor{#2}{\bfseries(}{\normalfont\sffamily\textit{##3}}\textcolor{#2}{\bfseries)}}}%
}
\NewIconThmStyle{purplethm}{EnvTheorem}{\faGem}
\NewIconThmStyle{defthm}{EnvTheorem}{\faBookOpen}
\NewIconThmStyle{examplethm}{EnvTheorem}{\faLightbulb[regular]}
\NewIconThmStyle{remarkthm}{EnvRemark}{\faInfoCircle}
\NewIconThmStyle{importantremarkthm}{EnvImportantRemark}{\faExclamationTriangle}
\NewIconThmStyle{questionthm}{EnvQuestion}{\faQuestionCircle}
\NewIconThmStyle{answerthm}{EnvQuestion}{\faCommentDots}

\theoremstyle{purplethm}
\newtheorem{theorem}{Satz}[chapter]
\newtheorem{lemma}[theorem]{Lemma}
\newtheorem{corollary}[theorem]{Korollar}
\newtheorem{proposition}[theorem]{Proposition}
\newtheorem{property}[theorem]{Eigenschaft}
\newtheorem{principle}[theorem]{Prinzip}
\newtheorem{fact}[theorem]{Fakt}
\newtheorem{claim}[theorem]{Behauptung}
\newtheorem*{theorem*}{Satz}
\newtheorem*{lemma*}{Lemma}
\newtheorem*{proposition*}{Proposition}
\newtheorem*{corollary*}{Korollar}
\newtheorem*{property*}{Eigenschaft}
\newtheorem*{principle*}{Prinzip}
\newtheorem*{fact*}{Fakt}
\newtheorem*{claim*}{Behauptung}

\theoremstyle{defthm}
\newtheorem{definition}[theorem]{Definition}
\newtheorem{axiom}[theorem]{Axiom}
\newtheorem{axioms}[theorem]{Axiome}
\newtheorem*{definition*}{Definition}
\newtheorem*{axiom*}{Axiom}
\newtheorem*{axioms*}{Axiome}

\theoremstyle{examplethm}
\newtheorem{example}[theorem]{Beispiel}
\newtheorem{examples}[theorem]{Beispiele}
\newtheorem{exercise}[theorem]{{\"U}bung}
\newtheorem{warmup}[theorem]{Einstieg}
\newtheorem*{example*}{Beispiel}
\newtheorem*{examples*}{Beispiele}
\newtheorem*{exercise*}{{\"U}bung}
\newtheorem*{warmup*}{Einstieg}

\theoremstyle{remarkthm}
\newtheorem{remark}[theorem]{Bemerkung}
\newtheorem{observation}[theorem]{Beobachtung}
\newtheorem{notation}[theorem]{Notation}
\newtheorem{summary}[theorem]{Zusammenfassung}
\newtheorem{goals}[theorem]{Ziele}
\newtheorem*{remark*}{Bemerkung}
\newtheorem*{observation*}{Beobachtung}
\newtheorem*{notation*}{Notation}
\newtheorem*{summary*}{Zusammenfassung}
\newtheorem*{goals*}{Ziele}

\theoremstyle{questionthm}
\newtheorem{question}[theorem]{Frage}
\theoremstyle{answerthm}
\newtheorem{answer}[theorem]{Antwort}
\theoremstyle{questionthm}
\newtheorem*{question*}{Frage}
\theoremstyle{answerthm}
\newtheorem*{answer*}{Antwort}

\theoremstyle{importantremarkthm}
\newtheorem{importantremark}[theorem]{Wichtige Bemerkung}
\newtheorem*{importantremark*}{Wichtige Bemerkung}

\theoremstyle{purplethm}

% Cleveref Namen (Deutsch)
\crefname{theorem}{Satz}{S{\"a}tze}
\Crefname{theorem}{Satz}{S{\"a}tze}
\crefname{lemma}{Lemma}{Lemmata}
\Crefname{lemma}{Lemma}{Lemmata}
\crefname{corollary}{Korollar}{Korollare}
\Crefname{corollary}{Korollar}{Korollare}
\crefname{proposition}{Proposition}{Propositionen}
\Crefname{proposition}{Proposition}{Propositionen}
\crefname{definition}{Definition}{Definitionen}
\Crefname{definition}{Definition}{Definitionen}
\crefname{example}{Beispiel}{Beispiele}
\Crefname{example}{Beispiel}{Beispiele}
\crefname{axiom}{Axiom}{Axiome}
\Crefname{axiom}{Axiom}{Axiome}
\crefname{axioms}{Axiome}{Axiome}
\Crefname{axioms}{Axiome}{Axiome}
\crefname{property}{Eigenschaft}{Eigenschaften}
\Crefname{property}{Eigenschaft}{Eigenschaften}
\crefname{principle}{Prinzip}{Prinzipien}
\Crefname{principle}{Prinzip}{Prinzipien}
\crefname{fact}{Fakt}{Fakten}
\Crefname{fact}{Fakt}{Fakten}
\crefname{claim}{Behauptung}{Behauptungen}
\Crefname{claim}{Behauptung}{Behauptungen}
\crefname{remark}{Bemerkung}{Bemerkungen}
\Crefname{remark}{Bemerkung}{Bemerkungen}
\crefname{observation}{Beobachtung}{Beobachtungen}
\Crefname{observation}{Beobachtung}{Beobachtungen}
\crefname{notation}{Notation}{Notationen}
\Crefname{notation}{Notation}{Notationen}
\crefname{summary}{Zusammenfassung}{Zusammenfassungen}
\Crefname{summary}{Zusammenfassung}{Zusammenfassungen}
\crefname{warmup}{Einstieg}{Einstiege}
\Crefname{warmup}{Einstieg}{Einstiege}
\crefname{exercise}{{\"U}bung}{{\"U}bungen}
\Crefname{exercise}{{\"U}bung}{{\"U}bungen}
\crefname{question}{Frage}{Fragen}
\Crefname{question}{Frage}{Fragen}
\crefname{answer}{Antwort}{Antworten}
\Crefname{answer}{Antwort}{Antworten}
\crefname{importantremark}{Wichtige Bemerkung}{Wichtige Bemerkungen}
\Crefname{importantremark}{Wichtige Bemerkung}{Wichtige Bemerkungen}
\crefname{goals}{Ziel}{Ziele}
\Crefname{goals}{Ziel}{Ziele}

% ==========================================
% VIBRANT SHADOW META-ENVIRONMENTS (Deutsch)
% ==========================================

% 1. Clean Spoken Protocol
\newtcolorbox{speech}[1][]{
enhanced,
frame hidden,
colback=EnvSpokenBg,
dynamic shadow,
code={
  \ifinsideproof
    \tcbset{borderline west={3pt}{0pt}{EnvSpokenProof}}
  \else
    \tcbset{borderline west={3pt}{0pt}{EnvSpoken}}
  \fi
},
title={\sffamily\bfseries\textcolor{\ifinsideproof EnvSpokenProof\else EnvSpoken\fi}{\faMicrophone\ \ Gesprochenes Wort\ifblank{#1}{}{ [#1]}}},
attach title to upper={\par\setlength{\parindent}{\standardparindent}\setlength{\parskip}{\customparskip}\vspace{1ex}},
breakable
}

% 2. Overview / Didactic Summary
\newtcolorbox{overview}[1][]{
enhanced,
colback=white,
boxrule=0.5pt,
dynamic shadow,
varwidth boxed title*=-30pt,
code={
  \ifinsideproof
    \tcbset{
      colframe=EnvNoteProof,
      boxed title style={colback=EnvNoteBg, colframe=EnvNoteProof, boxrule=0.5pt}
    }
  \else
    \tcbset{
      colframe=EnvNote,
      boxed title style={colback=EnvNoteBg, colframe=EnvNote, boxrule=0.5pt}
    }
  \fi
},
title={\sffamily\bfseries\begin{tabular}[t]{@{}l@{\hspace{0.5em}}l@{}}\faStickyNote & \begin{tabular}[t]{@{}l@{}} \ifblank{#1}{{\"U}berblick}{#1} \end{tabular}\end{tabular}},
coltitle=black,
attach boxed title to top left={xshift=10pt, yshift*=-\tcboxedtitleheight/2},
top=16pt,
breakable
}

% 3. Best Teacher / Didactic Insight
\newtcolorbox{insight}[1][Didaktischer Hinweis]{
enhanced,
arc=3mm,
frame hidden,
colback=EnvTeacherBg,
dynamic shadow,
code={
  \ifinsideproof
    \tcbset{borderline west={2pt}{0pt}{EnvTeacherProof}}
  \else
    \tcbset{borderline west={2pt}{0pt}{EnvTeacherBorder}}
  \fi
},
title={\sffamily\bfseries\textcolor{\ifinsideproof EnvTeacherProof\else EnvTeacherBorder\fi}{\textit{\faLightbulb\ \ Didaktischer Hinweis:} #1}},
attach title to upper=\par\vspace{0.5ex},
breakable
}

% 4a. Modular Visual Blackboard Replication Box
\NewTColorBox{highlight}{ m o }{
enhanced,
colback=white,
colupper=black,
fonttitle=\bfseries\sffamily,
coltitle=black,
boxrule=1.5pt,
arc=4pt,
dynamic color shadow=#1!30,
varwidth boxed title*=-30pt,
code={
  \ifinsideproof
    \tcbset{
      colframe=#1!50!ProofTint,
      boxed title style={colback=white, colframe=#1!50!ProofTint, boxrule=0.5pt, top=2pt, bottom=2pt, left=4pt, right=4pt}
    }
  \else
    \tcbset{
      colframe=#1,
      boxed title style={colback=white, colframe=#1, boxrule=0.5pt, top=2pt, bottom=2pt, left=4pt, right=4pt}
    }
  \fi
},
IfValueTF={#2}{
  title={\settowidth{\dimen0}{\faBox\ \ }\hangindent=\dimen0\hangafter=1\faBox\ \ #2},
  attach boxed title to top left={xshift=10pt, yshift=-10pt},
  % the attached title reaches 10pt into the box: keep the first line clear of it (as overview)
  top=14pt
}{
  notitle
}
}

% 4b. Legacy Visual Blackboard Replication Boxes
\newcommand{\NewChalkBox}[3]{%
\newtcolorbox{#1}[1][Hervorgehobene Tafelanschrift]{%
enhanced,
colback=white,
fonttitle=\bfseries\sffamily,
coltitle=black,
title={\settowidth{\dimen0}{\faBox\ \ }\hangindent=\dimen0\hangafter=1\faBox\ \ ##1},
boxrule=1.5pt,
arc=4pt,
dynamic color shadow=#3,
varwidth boxed title*=-30pt,
attach boxed title to top left={xshift=10pt, yshift=-10pt},
code={
  \ifinsideproof
    \tcbset{
      colframe=#2!50!ProofTint,
      boxed title style={colback=white, colframe=#2!50!ProofTint, boxrule=0.5pt, top=2pt, bottom=2pt, left=4pt, right=4pt}
    }
  \else
    \tcbset{
      colframe=#2,
      boxed title style={colback=white, colframe=#2, boxrule=0.5pt, top=2pt, bottom=2pt, left=4pt, right=4pt}
    }
  \fi
}
}%
}

\NewChalkBox{orange-box}{BurntOrange}{BurntOrange!30}
\NewChalkBox{blue-box}{MidnightBlue}{MidnightBlue!30}
\NewChalkBox{dark-green-box}{ForestGreen}{ForestGreen!30}
\NewChalkBox{apple-green-box}{LimeGreen}{LimeGreen!30}
\NewChalkBox{yellow-box}{Goldenrod!90!black}{Goldenrod!30}
\NewChalkBox{violet-box}{Violet}{Violet!30}
\NewChalkBox{red-box}{BrickRed}{BrickRed!30}
\NewChalkBox{green-box}{Green}{Green!30}
\NewChalkBox{purple-box}{Purple}{Purple!30}
\NewChalkBox{gray-box}{Gray}{Gray!30}
\NewChalkBox{apricot-box}{Apricot}{Apricot!30}
\NewChalkBox{bittersweet-box}{Bittersweet}{Bittersweet!30}
\NewChalkBox{teal-blue-box}{TealBlue}{TealBlue!30}
\NewChalkBox{cyan-box}{Cyan}{Cyan!30}
\NewChalkBox{magenta-box}{Magenta}{Magenta!30}
\NewChalkBox{brown-box}{Brown}{Brown!30}
\NewChalkBox{pink-box}{Pink}{Pink!30}

\newcommand{\NewFormulaBox}[3]{%
\newtcolorbox{#1}[1][Hervorgehobene Formel]{%
enhanced,
colback=white,
fonttitle=\bfseries\sffamily,
coltitle=black,
title={\settowidth{\dimen0}{\faIcon{square-root-alt}\ \ }\hangindent=\dimen0\hangafter=1\faIcon{square-root-alt}\ \ ##1},
boxrule=1.5pt,
arc=4pt,
dynamic color shadow=#3,
varwidth boxed title*=-30pt,
attach boxed title to top left={xshift=10pt, yshift=-10pt},
code={
  \ifinsideproof
    \tcbset{
      colframe=#2!50!ProofTint,
      boxed title style={colback=white, colframe=#2!50!ProofTint, boxrule=0.5pt, top=2pt, bottom=2pt, left=4pt, right=4pt}
    }
  \else
    \tcbset{
      colframe=#2,
      boxed title style={colback=white, colframe=#2, boxrule=0.5pt, top=2pt, bottom=2pt, left=4pt, right=4pt}
    }
  \fi
}
}%
}

\NewFormulaBox{orange-formula}{BurntOrange}{BurntOrange!30}
\NewFormulaBox{blue-formula}{MidnightBlue}{MidnightBlue!30}
\NewFormulaBox{dark-green-formula}{ForestGreen}{ForestGreen!30}
\NewFormulaBox{apple-green-formula}{LimeGreen}{LimeGreen!30}
\NewFormulaBox{yellow-formula}{Goldenrod!90!black}{Goldenrod!30}
\NewFormulaBox{violet-formula}{Violet}{Violet!30}
\NewFormulaBox{red-formula}{BrickRed}{BrickRed!30}

% 5. Mathematical Setup (Stroke-by-Stroke)
\NewTColorBox{content}{ O{} }{
enhanced,
frame hidden,
colback=white,
opacityback=0, 
breakable,
before upper={\setlength{\parskip}{\customparskip}\setlength{\parindent}{0pt}},
code={
  \ifinsideproof
    \tcbset{borderline west={2pt}{0pt}{EnvMathProof}}
  \else
    \tcbset{borderline west={2pt}{0pt}{EnvMath!50}}
  \fi
  \ifblank{#1}{
    \tcbset{notitle}
  }{
    \tcbset{title={\sffamily\bfseries\textcolor{\ifinsideproof EnvMathProof\else EnvMath\fi}{\faChalkboard\ \ #1}}, 
    attach title to upper={\par\setlength{\parindent}{0pt}\setlength{\parskip}{\customparskip}\vspace{1ex}}}
  }
}
}

% 6. Meta-Note
\newtcolorbox{note}[1][Organisatorische Metanotiz]{
enhanced,
frame hidden,
colback=white,
opacityback=0,
code={
  \ifinsideproof
    \tcbset{borderline={0.5pt}{0pt}{EnvMetaProof, dashed}}
  \else
    \tcbset{borderline={0.5pt}{0pt}{EnvMeta, dashed}}
  \fi
},
title={\sffamily\bfseries\textcolor{\ifinsideproof EnvMetaProof\else EnvMeta\fi}{\faEye\ \ #1}},
attach title to upper=\par\vspace{1ex},
breakable
}

% 7. AI Observational Note
\newtcolorbox{aiNote}[1][KI-Notiz]{
enhanced,
colback=EnvAIBg,
dynamic shadow,
code={
  \ifinsideproof
    \tcbset{colframe=EnvAIProof}
  \else
    \tcbset{colframe=EnvAIBorder}
  \fi
},
fonttitle=\bfseries\sffamily,
coltitle=white,
title={\faRobot\ \ #1},
boxrule=1pt,
arc=3pt,
breakable
}

% 8. Student Question
\newtcolorbox{student}[1][Frage aus dem Publikum]{
enhanced,
frame hidden,
colback=EnvStudentBg,
dynamic shadow,
code={
  \ifinsideproof
    \tcbset{borderline west={3pt}{0pt}{EnvStudentProof}}
  \else
    \tcbset{borderline west={3pt}{0pt}{EnvStudent}}
  \fi
},
title={\sffamily\bfseries\textcolor{\ifinsideproof EnvStudentProof\else EnvStudent\fi}{\faUserGraduate\ \ #1}},
attach title to upper=\par\vspace{1ex},
breakable
}

% 9. Explanation of Steps
\newenvironment{steps}[1][Erkl{\"a}rung der Schritte]{
\par\vspace{1ex}\noindent\textbf{\textcolor{ExplanationIconColor}{\faListUl\ \ #1:}}
}{
\par\vspace{1ex}
}

% 10. Subsection-like heading for proofs
\newenvironment{proof-subsection}[1]{
  \par\vspace{1.5ex}\noindent\textbf{\textcolor{ProofTint}{\faAngleRight\ \ #1}}\par\vspace{0.5ex}
}{
  \par\vspace{1.5ex}
}

% 11. Proofs and QED
\makeatletter
\newcommand{\qedbadge}{%
  \tikz[baseline=(Q.base)]\node[draw=ProofTint, fill=ProofQEDBg, line width=1pt, rounded corners=2pt, inner xsep=4pt, inner ysep=2pt, font=\sffamily\bfseries\small, text=ProofQEDText, drop shadow={shadow xshift=0.5ex, shadow yshift=-0.5ex, fill=ProofTint, opacity=0.4}] (Q) {Q.E.D.};%
}

% SHORT PROOF 
\renewenvironment{proof}[1][Beweis]{%
    \par\addvspace{2ex plus 0.5ex minus 0.2ex}%
    \insideprooftrue 
    \noindent\textcolor{ProofTint}{\faPenNib}\hspace{0.5em}%
    {\normalfont\sffamily\bfseries\color{ProofTint}#1}%
    \par\nobreak\vspace{1ex plus .2ex}\vspace{-\parskip}%
    \@afterindentfalse\@afterheading
}{%
    \ifmmode
      \mathqed
    \else
      \ifhmode
        \unskip\nobreak\hfill\penalty50\hskip2em\hbox{}\nobreak\hfill\qedbadge\par\vspace{2ex}%
      \else
        \par\vspace{-\parskip}\noindent\hfill\qedbadge\par\vspace{2ex}%
      \fi
    \fi
}

% LONG PROOF
\newcommand*{\proofWrapperDefaultTitle}{Beweis}
\let\proofwrapper@tocclose\relax
\newcommand{\proofwrapper@depth}{%
  \ifnum\value{section}=\z@ 0%
  \else\ifnum\value{subsection}=\z@ 1%
  \else\ifnum\value{subsubsection}=\z@ 2%
  \else 3\fi\fi\fi}
\newcommand{\proofwrapper@tocentry}[1]{%
  \edef\proofwrapper@level{\proofwrapper@depth}%
  \phantomsection
  \addtocontents{toc}{\protect\proofwrappertocbegin{\proofwrapper@level}}%
  \edef\toclevel@proofwrapper{\the\numexpr\proofwrapper@level+1\relax}%
  \addcontentsline{toc}{proofwrapper}{\protect\proofwrappertoctitle{#1}}%
  % Bookmarks of the headings inside the proof become children of the proof's bookmark, by the
  % same rule as in the TOC: the first heading fixes the shift. The bookmark depth limit moves
  % with them, so that exactly the levels shown outside a proof are shown inside it.
  \gdef\proofwrapper@bmshift{-1}%
  \let\proofwrapper@addcontentsline\addcontentsline
  \def\addcontentsline##1##2##3{\proofwrapper@bmlevel{##2}\proofwrapper@addcontentsline{##1}{##2}{##3}}%
  \def\proofwrapper@tocclose{\addtocontents{toc}{\protect\proofwrappertocend}}%
}
\def\proofwrapper@baselevel@section{1}
\def\proofwrapper@baselevel@subsection{2}
\def\proofwrapper@baselevel@subsubsection{3}
\def\proofwrapper@baselevel@paragraph{4}
\def\proofwrapper@baselevel@subparagraph{5}
\newcommand{\proofwrapper@bmlevel}[1]{%
  \ifcsname proofwrapper@baselevel@#1\endcsname
    \ifnum\proofwrapper@bmshift<\z@
      \count@=\numexpr\proofwrapper@level+2-\csname proofwrapper@baselevel@#1\endcsname\relax
      \ifnum\count@<\z@ \count@=\z@ \fi
      \xdef\proofwrapper@bmshift{\the\count@}%
    \fi
    \expandafter\edef\csname toclevel@#1\endcsname
      {\the\numexpr\csname proofwrapper@baselevel@#1\endcsname+\proofwrapper@bmshift\relax}%
    \edef\Hy@bookmarksdepth{\the\numexpr\c@tocdepth+\proofwrapper@bmshift\relax}%
  \fi}
% Inside the TOC entry, references are printed without a link of their own: the entry is a link.
\newcommand{\proofwrappertoctitle}[1]{{\let\proofwrapper@cref\cref\let\proofwrapper@Cref\Cref
  \def\cref{\proofwrapper@cref*}\def\Cref{\proofwrapper@Cref*}#1}}
% PDF bookmarks cannot hold a \cref; spell it out from the label data cleveref writes to the aux
% file ([theorem][4][3]3.4 -> "Satz 3.4"), or "??" before the first run has written it.
\newcommand{\proofwrapper@pdfref}[1]{%
  \ifcsname r@#1@cref\endcsname
    \expandafter\@firstoftwo
  \else
    \expandafter\@secondoftwo
  \fi
  {\expandafter\expandafter\expandafter\proofwrapper@pdfref@i\csname r@#1@cref\endcsname\@nil}%
  {??}}
\def\proofwrapper@pdfref@i#1#2\@nil{\proofwrapper@pdfref@ii#1\@nil}
\def\proofwrapper@pdfref@ii[#1][#2][#3]#4\@nil{\csname Cref@#1@name\endcsname\space#4}
\pdfstringdefDisableCommands{%
  \let\proofwrappertoctitle\@firstofone
  \let\cref\proofwrapper@pdfref
  \let\Cref\proofwrapper@pdfref
}

\newenvironment{proofWrapper}[1][Beweis]{%
    \par\addvspace{3.5ex plus 1ex minus .2ex}%
    \insideprooftrue
    \noindent
    \def\proofwrapper@title{#1}%
    \ifx\proofwrapper@title\proofWrapperDefaultTitle\else\proofwrapper@tocentry{#1}\fi
    \makebox[0pt][r]{\textcolor{ProofTint}{\faPenNib}\hspace{10pt}}%
    {\normalfont\sffamily\Large\bfseries\color{ProofTint}#1}%
    \par\nobreak\vspace{1.5ex plus .2ex}\vspace{-\parskip}%
    \@afterindentfalse\@afterheading
}{%
    \ifmmode
      \mathqed
    \else
      \ifhmode
        \unskip\nobreak\hfill\penalty50\hskip2em\hbox{}\nobreak\hfill\qedbadge\par\vspace{2ex}%
      \else
        \par\vspace{-\parskip}\noindent\hfill\qedbadge\par\vspace{2ex}%
      \fi
    \fi
    \proofwrapper@tocclose
}
\makeatother

% 13. Lecture Break
\newtcolorbox{pause}[1][Vorlesungspause]{
enhanced,
frame hidden,
colback=white,
opacityback=0,
code={
  \ifinsideproof
    \tcbset{borderline north={0.5pt}{0pt}{EnvBreakProof, dashed}, borderline south={0.5pt}{0pt}{EnvBreakProof, dashed}}
  \else
    \tcbset{borderline north={0.5pt}{0pt}{EnvBreak, dashed}, borderline south={0.5pt}{0pt}{EnvBreak, dashed}}
  \fi
},
title={\sffamily\bfseries\centering\textcolor{\ifinsideproof EnvBreakProof\else EnvBreak\fi}{\faCoffee\ \ #1}},
attach title to upper=\par\vspace{0.5ex},
fontupper=\centering\itshape\color{\ifinsideproof EnvBreakProof\else EnvBreak\fi},
breakable
}

% 14. Fallback Boxes
\newcommand{\NewFallbackBox}[1]{%
\newtcolorbox{#1}[1][Warnung]{
enhanced,
colback=EnvAIBg,
dynamic shadow,
code={
  \ifinsideproof
    \tcbset{colframe=EnvAIFallbackProof}
  \else
    \tcbset{colframe=EnvAI}
  \fi
},
fonttitle=\sffamily\bfseries,
coltitle=white,
title={\faExclamationTriangle\ \ ##1},
breakable
}
}
\NewFallbackBox{ai-logical-gap}
\NewFallbackBox{ai-generation-loop-fallback}
\NewFallbackBox{ai-raw-ocr-fallback}
\NewFallbackBox{ai-phonetic-dump}
\NewFallbackBox{ai-modality-conflict}
\NewFallbackBox{ai-off-camera-state}
\NewFallbackBox{ai-async-board-update}
\NewFallbackBox{ai-kinetic-emphasis}
\NewEnviron{hidden}{}

% Custom Lecture Chapter Command with Optional Semester/Week Number
\newif\ifrestorelecturesubcounters
\newcounter{currentweek}
\setcounter{currentweek}{0}
\newcounter{lecturecount}
\setcounter{lecturecount}{0}
\renewcommand{\theHchapter}{\thechapter.\arabic{lecturecount}}

\NewDocumentCommand{\lecturechapter}{ o m m m m }{%
  \restorelecturesubcountersfalse
  \stepcounter{lecturecount}%
  \IfValueT{#1}{%
    \ifblank{#1}{}{%
      \ifnum\numexpr#1\relax>\value{currentweek}%
        \setcounter{currentweek}{#1}%
        \phantomsection
        \addcontentsline{toc}{part}{Woche #1}%
      \fi
      \ifnum\numexpr#1\relax=\value{chapter}%
        \restorelecturesubcounterstrue
        \edef\TempSec{\the\value{section}}%
        \edef\TempThm{\the\value{theorem}}%
        \edef\TempEq{\the\value{equation}}%
        \edef\TempFig{\the\value{figure}}%
        \edef\TempTab{\the\value{table}}%
      \fi
      \setcounter{chapter}{\numexpr#1-1\relax}%
    }%
  }%
  \chapter[\texorpdfstring{#2, #3\quad---\quad #5}{#2, #3 --- #5}]{{\Large\color{ChapSubtitleText}#2 #4} \texorpdfstring{\\[1ex]}{--- } {\LARGE #5}}%
  \ifrestorelecturesubcounters
    \setcounter{section}{\TempSec}%
    \setcounter{theorem}{\TempThm}%
    \setcounter{equation}{\TempEq}%
    \setcounter{figure}{\TempFig}%
    \setcounter{table}{\TempTab}%
  \fi
  \markboth{\chaptername\ \thechapter\ \ ·\ \ #2, #3}{}% the title itself is on the chapter page
  \InsertMark{gs-sec}{}\InsertMark{gs-subsec}{}\InsertMark{gs-subsubsec}{}%
}

% YouTube / standalone video chapter: \youtubechapter[Date]{Creator / Channel}{Video Title}
\NewDocumentCommand{\youtubechapter}{ o m m }{%
  \stepcounter{lecturecount}%
  \IfValueTF{#1}{%
    \chapter[\texorpdfstring{#2, #1\quad---\quad #3}{#2, #1 --- #3}]{{\Large\color{ChapSubtitleText}#2\ ---\ #1} \texorpdfstring{\\[1ex]}{--- } {\LARGE #3}}%
    \markboth{\chaptername\ \thechapter\ \ ·\ \ #2, #1}{}%
  }{%
    \chapter[\texorpdfstring{#2\quad---\quad #3}{#2 --- #3}]{{\Large\color{ChapSubtitleText}#2} \texorpdfstring{\\[1ex]}{--- } {\LARGE #3}}%
    \markboth{\chaptername\ \thechapter\ \ ·\ \ #2}{}%
  }%
  \InsertMark{gs-sec}{}\InsertMark{gs-subsec}{}\InsertMark{gs-subsubsec}{}%
}

\newcommand{\chapterlabel}{\MakeUppercase{\chaptertitlename}\ \thechapter}
\titlespacing*{\chapter}{0pt}{-20pt}{18pt plus 4pt minus 2pt}

\titleformat{\chapter}[display]
{\normalfont\sffamily\Large\bfseries\color{ChapTitleText}}
{\raggedleft\large\color{ChapNumberText}\chapterlabel}
{1ex}
{\color{ChapLine}\rule{\textwidth}{1.5pt}\endgraf\nopagebreak\vspace{2ex}\raggedleft} 
[\endgraf\nopagebreak\vspace{1ex}\color{ChapLine}\rule{\textwidth}{1.5pt}]

% Extended Section Bars 
\newcommand{\sectionbar}{%
\tikz[baseline=0pt]{
\fill[\ifinsideproof SecBar1Proof\else SecBar1\fi] (-2.5pt, 0) rectangle (1.0pt, 0.72em);
\fill[\ifinsideproof SecBar2Proof\else SecBar2\fi] (2.0pt, 0) rectangle (3.5pt, 0.72em);
\fill[\ifinsideproof SecBar3Proof\else SecBar3\fi] (4.8pt, 0) rectangle (5.3pt, 0.72em);
\fill[\ifinsideproof SecBar4Proof\else SecBar4\fi] (6.6pt, 0) rectangle (8.1pt, 0.72em);
}%
\hspace{10pt}%
}

\newcommand{\subsectionbar}{%
\tikz[baseline=0pt]{
\fill[\ifinsideproof SubSecBar1Proof\else SubSecBar1\fi] (-2.5pt, 0) rectangle (0.3pt, 0.72em);
\shade[left color=\ifinsideproof SubSecBar1Proof\else SubSecBar1\fi, right color=\ifinsideproof SubSecBar2Proof\else SubSecBar2\fi] (1.3pt, 0) rectangle (2.0pt, 0.72em);
\fill[\ifinsideproof SubSecBar2Proof\else SubSecBar2\fi] (3.3pt, 0) rectangle (4.8pt, 0.72em);
}%
\hspace{10pt}%
}

\newcommand{\subsubsectionbar}{%
\tikz[baseline=0pt]{
\fill[\ifinsideproof SubSubSecBar1Proof\else SubSubSecBar1\fi] (-2.5pt, 0) rectangle (-1.0pt, 0.72em);
\fill[\ifinsideproof SubSubSecBar2Proof\else SubSubSecBar2\fi] (0,0) rectangle (2.2pt, 0.72em);
}%
\hspace{10pt}%
}

% Headings show no number, as in the other books; the numbers stay in the TOC, the running
% heads and \cref. The \color sits in the before-code (horizontal mode), not in the format
% argument: there it would add a color whatsit to the vertical list, and the glue after it is a
% legal page break, which let a \section heading be separated from the \subsection below it.
\titleformat{\section}{\normalfont\sffamily\Large\bfseries}{}{0pt}{\phantomsection\color{\ifinsideproof SecTitleColorProof\else SecTitleColor\fi}\makebox[0pt][r]{\sectionbar}}
\titleformat{\subsection}{\normalfont\sffamily\large\bfseries}{}{0pt}{\phantomsection\color{\ifinsideproof SubSecTitleColorProof\else SubSecTitleColor\fi}\makebox[0pt][r]{\subsectionbar}}
\titleformat{\subsubsection}{\normalfont\sffamily\normalsize\bfseries}{}{0pt}{\phantomsection\color{\ifinsideproof SubSubSecTitleColorProof\else SubSubSecTitleColor\fi}\makebox[0pt][r]{\subsubsectionbar}}
\titleformat{\paragraph}[runin]{\normalfont\normalsize\bfseries}{\theparagraph}{1em}{\phantomsection} % the class default, plus a link target on the heading (unnumbered, it has none)

\titlespacing*{\section}{0pt}{3.5ex plus 1ex minus .2ex}{1.5ex plus .2ex}
\titlespacing*{\subsection}{0pt}{3ex plus 1ex minus .2ex}{1.2ex plus .2ex}
\titlespacing*{\subsubsection}{0pt}{2.5ex plus 1ex minus .2ex}{1ex plus .2ex}

% =====================================================================
% BACKWARD-COMPATIBILITY ALIASES
% =====================================================================
\newcommand{\AliasEnv}[2]{%
  \newenvironment{#1}{\begin{#2}}{\end{#2}}%
  \newenvironment{#1*}{\begin{#2*}}{\end{#2*}}%
}

\AliasEnv{spoken-clean}{speech}
\AliasEnv{math-stroke}{content}
\AliasEnv{student-interaction}{student}
\AliasEnv{explanation-of-steps}{steps}
\AliasEnv{didactic-insight}{insight}
\AliasEnv{meta-note}{note}
\AliasEnv{nice-box}{overview}
\AliasEnv{color-box}{highlight}
\AliasEnv{lecture-break}{pause}
\AliasEnv{ai-note}{aiNote}
\AliasEnv{short-proof}{proof}

\NewEnviron{ai-tikz-planner-invisible-content}{}
\NewEnviron{ai-proof-skeleton-invisible-content}{}
\NewEnviron{ai-global-state-checkpoint-invisible-content}{}
\NewEnviron{ai-retroactive-patch}[1][]{}
\NewEnviron{ai-global-state-checkpoint}{}
